\section{A smooth projective application}\label{sec:descent}

A section on a compact factor of the universal cover descends only if it respects the deck action. The invariant conclusion of Theorem~\ref{thm:equivariant} supplies precisely this compatibility with the residual torus monodromy. We obtain another proof of the smooth nonvanishing result established in \cite[Theorem ``Anticanonical nonvanishing'']{SmoothCompanion} by a finite-volume argument. Here we use that article's finite-cover structure theorem and norm lemma; the invariant section comes from the bounded-metric criterion and stable quotient constructed above.

\begin{corollary}\label{cor:smooth}
Let $X$ be a smooth connected projective complex variety with smoothly semipositive anticanonical bundle. Then $H^0(X,-mK_X)\ne0$ for some integer $m>0$.
\end{corollary}

\begin{proof}
Use the finite-cover structure theorem of \cite[Theorem ``Finite cover with compact torus monodromy'']{SmoothCompanion}, in its form retaining the deck action. There is a finite \'etale cover $X'\to X$ whose universal cover is
\[
\widetilde X=\CC^q\times B\times Z.
\]
The compact factors are projective; $K_B$ is trivial, and $Z$ is rationally connected with smoothly semipositive anticanonical bundle. Either factor may be a point. The deck group acts by translations on $\CC^q$, trivially on $B$, and on $Z$ through a compact torus contained in an algebraic torus $T$. The Euclidean and $B$ factors have invariant nowhere zero canonical frames.

By Theorem~\ref{thm:equivariant}, choose a nonzero $T$-invariant section $s_Z$ of $-mK_Z$ for some $m>0$. If $Z$ is a point, use the constant section. Tensor $s_Z$ with the $m$th powers of the dual canonical frames on the other two factors. This yields a nonzero section of $-mK_{\widetilde X}$ invariant under the deck group of $\widetilde X\to X'$. Natural linearization is crucial here: the action on the anticanonical product is the differential action used in Theorem~\ref{thm:equivariant}. The section descends to $X'$, and projectivity makes it algebraic. If $d$ is the degree of $X'\to X$, the finite \'etale norm \cite[Lemma ``Finite \'etale norm'']{SmoothCompanion} gives a nonzero section of $-mdK_X$.
\end{proof}

The geometric input in this proof rests on Yau's prescribed-Ricci theorem \cite{Yau1978} and the structure results of Demailly--Peternell--Schneider \cite{DPS1996} and Campana--Demailly--Peternell \cite[Theorem~1.4]{CDP2015}. A finite cover need not be a product: the residual torus monodromy is the reason that an invariant section, rather than an arbitrary section on $Z$, is needed.

The positive multiple in Corollary~\ref{cor:smooth} cannot in general be replaced by the first power. For an Enriques surface $S$, $K_S$ is a nontrivial line bundle of order two and $H^0(S,-K_S)=0$ \cite{Dolgachev2016}. A flat metric on $-K_S$ and the Fubini--Study metric on $\OO_{\mathbb P^1}(2)$ make the anticanonical bundle of $S\times\mathbb P^1$ smoothly semipositive. The product formula for sections shows that it has no first-power section, while its second power has sections. Its restriction to a $\mathbb P^1$ fiber has degree two, so this example is not explained by torsion of the anticanonical bundle itself.
